\begin{lemma}[Converse game]
Let $Q$ be a type, let $r$ be a relation on $Q$, and let $F$ be a front on
$\mathbb N$ with $\emptyset\notin F$.  Let
$f:F\to Q$ be a super-sequence, and suppose that $f$ is bad.  Define the
hereditarily countable singleton sequence
\[
a(k)=\left\langle\widetilde f(\{k\}),\
  \mathsf{HereditarilyCountable}(\widetilde f(\{k\}))\right\rangle
  \in\mathsf{HerCtblPower}(Q),
\]
where the hereditary witness is supplied by
\noderef{TildeFHerCtbl} and singleton membership by
\noderef{FrontTreeSingleton}.  Then $a$ is bad for
$\mathsf{HerCtblRel}(r)$: for all $m<n$,
\[
\neg\mathsf{HerCtblRel}(r)(a(m),a(n)).
\]
\end{lemma}
\begin{proof}
Fix $m<n$ and suppose, toward a contradiction, that
\mathsf{HerCtblRel}(r)(a(m),a(n)).  By the definition of
\noderef{HerCtblRel}, this is the comparison
\[
\mathsf{PowerRel}(r)(\widetilde f(\{m\}),\widetilde f(\{n\})).
\]
The singleton terms are defined by \noderef{TildeFSingleton}, and their
hereditary witnesses are legitimate by \noderef{TildeFHerCtbl}.

We run the copying game from the paper.  Start with
$s_0=\{m\}$, $t_0=\{n\}$, and $u_0=\{m,n\}$.  By the structure of a
super-sequence \noderef{SuperSequence}, $f$ supplies a front on the full
base.  Thus \noderef{FrontBridgeInit} yields
\[
\mathsf{FrontBridge}(F,s_0,t_0,u_0).
\]
We maintain this bridge while preserving a comparison
\[
\mathsf{PowerRel}(r)(\widetilde f(s),\widetilde f(t))             \tag{*}
\]
for the current tree nodes $s,t$.  The induction is well-founded on the
lexicographic product of the two tree relations
\[
\mathsf{ProperInitialSegment}(s,s')
\quad\text{and}\quad
\mathsf{ProperInitialSegment}(t,t'),
\]
with a recursive call made whenever one coordinate is replaced by a
strict extension.  This is well-founded by
\noderef{FrontTreeWellFounded}; its product lexicographic relation is
therefore well-founded as well.

Suppose first that $s\notin F$.  The bridge definition
\noderef{FrontBridge} gives an ordered one-point extension $u$ of $s$,
with $u\in T(F)$.  Thus $u$ is one of the children in the non-front
recursion equation \noderef{TildeFNonfrontEquation}.  The same legality is
also the conclusion of \noderef{FrontBridgeStep}, and the all-extensions
form \noderef{FrontTreeOrderedExtension} ensures that no extra choice of
front member is being made.  If also $t\notin F$, the two values in (*)
are nodes.  By \noderef{PowerRelNodeNode}, the comparison supplies a right
child $t'$ with
\[
\mathsf{PowerRel}(r)(\widetilde f(u),\widetilde f(t')).
\]
Write this child as $t'=t\cup\{n\}$.  Since $t$ is the finite tail of
$u$ in the nonterminal bridge clause, the ordered-extension condition makes
$n$ larger than every member of $t$, hence larger than every member of $u$.
Thus this is exactly the first response alternative in
\noderef{FrontBridgeStep}.  Hence
\noderef{FrontBridgeStep} supplies an updated bridge
\mathsf{FrontBridge}(F,u,t',u')$, and the induction hypothesis applies
to $(u,t')$.  The left coordinate has strictly extended $s$.

If instead $t\in F$, the right value in (*) is an atom by
\noderef{TildeFFrontEquation}.  The node--atom reduction
\noderef{PowerRelNodeAtom} says that every left child is related to this
atom, in particular the prescribed child $u$.  We append a fresh point
above the current bridge and retain the terminal right node.  The second
case of \noderef{FrontBridgeStep} gives the resulting bridge, and the
induction hypothesis again applies to the strict left extension $(u,t)$.

Now suppose that $s\in F$ and $t\notin F$.  The left value is an atom by
\noderef{TildeFFrontEquation}.  By \noderef{PowerRelAtomNode}, the
comparison (*) supplies a child $t'$ of the right node with
\[
\mathsf{PowerRel}(r)(\widetilde f(s),\widetilde f(t')).
\]
Write the child as $t'=t\cup\{n\}$.  As above, the tail clause in
\noderef{FrontBridge} and the ordered-extension condition put $n$ above
every member of the bridge.  Choose it as the right response in
\noderef{FrontBridgeStep}, keep the terminal left node,
and obtain a bridge for $(s,t')$.  This is a strict extension in the
right coordinate, so the induction hypothesis applies.

It remains to consider $s,t\in F$.  The front equation
\noderef{TildeFFrontEquation} and the atom--atom reduction
\noderef{PowerRelAtomAtom} turn (*) into
\[
r\bigl(f(\langle s,\,s\in F\rangle),
       f(\langle t,\,t\in F\rangle)\bigr).
\]
The bridge has now reached two terminal front nodes.  By
\noderef{FrontBridgeTerminal}, it supplies
$\mathsf{FiniteShift}(s,t)$.  This contradicts the badness hypothesis
\noderef{BadSuperSequence}, applied to $s,t$, their displayed front
membership proofs, and this finite shift.  Thus the terminal case is
impossible, and well-founded induction rules out the original comparison.
Since $m<n$ was arbitrary, the sequence $a$ is bad as claimed.
\end{proof}
